% !TEX root = ../../tree_notes.tex
% =====================================================================
%  7.4  Classic Heap Problems                    sec:c07-classics
%  Exports: fig:c07-multimerge  lst:c07-multimerge  lst:c07-median
%           prob:c07-mergek  lst:c07-mergek-soln
%           (practices: LC 373, LC 378, LC 502)
%  Anti-repetition: the running median problem LC 295 is PLACED in Chapter 8's
%  mixed set -- cited by chapter, never re-listed here; the sliding-window MAX
%  is Chapter 9's monotonic deque (\chref{c26}) -- contrasted, not re-taught.
% =====================================================================
\section{Classic Heap Problems}\label{sec:c07-classics}

\begin{previouslybox}
\begin{enumerate}[nosep, leftmargin=1.6em]
  \item A merge of \emph{two} sorted lists walks both with two pointers. If you had $k$ lists, what would
        you re-examine at every output step? (the current head of each list -- $k$ candidates)
  \item \probref{c07-stone} and \probref{c07-connect} both ``pop the extreme, push something back.'' Which
        heap -- min or max -- matches ``smallest sum first''? (min-heap)
  \item To keep the top $k$ you used a size-$k$ heap (\secref{c07-build}). What is the \emph{dual}: keeping
        the $k$ \emph{smallest pair sums} out of two arrays' $n^2$ products without enumerating them?
        (best-first: a frontier min-heap seeded at the smallest cell)
  \item A median needs \emph{both} halves ordered. What two heaps split a stream into ``low half'' and
        ``high half''? (a max-heap for the low half, min-heap for the high -- built here)
\end{enumerate}
\end{previouslybox}

\begin{topicintro}
We now hold the tool -- a heap that hands back one extreme at a time -- and spend it on the four problem
families that recur for the rest of your competitive life. They are not four tricks; they are one idea
(keep the \emph{candidates for the next output} in a heap and extract) seen from four angles.

\textbf{Multiway merge.} Merging $k$ sorted sequences: at each step the next element is the smallest of the
$k$ current heads, so a \emph{min-heap of heads} turns a $\bigO{Nk}$ pairwise scan into $\bigO{N\log k}$ --
the anchor, LeetCode 23 (\probref{c07-mergek}), and the reason external sorting and $k$-way file merging
scale. \textbf{Best-first / $k$-best.} When the answer is the $k$ smallest objects of a huge implicit space
(pair sums, matrix cells), do not enumerate -- seed a frontier heap at the smallest cell and expand its
neighbours, popping $k$ times (\secref{c07-classics}'s practices LC 373, LC 378). \textbf{Greedy with
repair.} IPO-style problems where each choice unlocks new options: a max-heap of currently-affordable items,
refilled as your budget grows. \textbf{Two-heap median.} A streaming median is a max-heap for the low half
and a min-heap for the high half, balanced to within one -- the pattern we introduce here, though the
canonical problem (LeetCode 295) is placed in Chapter~8's mixed set (\chref{c06}) and is only cited.
\textbf{Prerequisites.} \secref{c07-heap} and \secref{c07-build}; the greedy-extraction reflex from
\probref{c07-stone}.

\textbf{After this section you can}
\begin{itemize}[nosep]
  \item merge $k$ sorted sequences in $\bigO{N\log k}$ with a head-heap, and adapt it to $k$-way external
        merge; \hfill\lv{4}
  \item answer ``$k$ smallest of an implicit grid/pair-set'' by best-first expansion, avoiding enumeration; \hfill\lv{4}--\lv{5}
  \item maintain a running median with two balanced heaps; \hfill\lv{4}
  \item recognise when a heap beats Chapter~9's monotonic deque (extreme of an \emph{editing} set vs.\ a
        \emph{sliding window}). \hfill\lv{4}
\end{itemize}
\end{topicintro}

\begin{recognitionbox}[Pattern card -- emit the next-best of many sources]
\reckey{you must \emph{output items in sorted / increasing order of some key}, but they arrive from several
pre-sorted streams, or live in an implicit space too large to enumerate, or you maintain a running
\emph{middle} rather than an extreme.}
\recuse{
  \item ``merge $k$ sorted lists / files'', ``smallest range touching all $k$ lists'' -- head-heap multiway
        merge, \probref{c07-mergek};
  \item ``$k$ pairs with smallest sum'', ``$k$-th smallest in a sorted matrix'' -- best-first frontier heap;
  \item ``median / running quantile of a stream'' -- two heaps; ``max profit given capital'' -- greedy with
        repair (IPO).
}
\recnotuse{
  \item the window is a fixed-size \emph{sliding} one and you want its max/min -- that is the monotonic
        \emph{deque} of Chapter~9 (\chref{c26}): elements leave by \emph{position}, not by value, so a heap's
        $\bigO{\log n}$ eviction of an arbitrary element is wasted;
  \item all $N$ items are known and memory is no issue -- a single $\bigO{N\log N}$ sort beats $k$-way merge;
  \item you need the top $k$ of \emph{one} stream by frequency -- that is \probref{c07-topk}'s size-$k$ heap,
        not a merge.
}
One-line test: \emph{are you producing the next-best item, repeatedly, from many places at once? A min-heap of
``the best each source can still offer'' is the merge; a min-heap of ``the frontier still unexplored'' is
best-first.}
\end{recognitionbox}

\subsection{Intuition: The Heap as an Output Scheduler}

Merging two sorted lists is a two-pointer race -- compare the two heads, emit the smaller. With $k$ lists a
naive scan re-compares all $k$ heads at every one of the $N$ output positions, costing $\bigO{Nk}$. But you
never actually need to \emph{see} all $k$ heads; you only need the smallest one, and then, once it is
emitted, the \emph{next} smallest among the (now updated) heads. ``Give me the current smallest of a set that
changes as I emit'' is precisely \code{top()}--\code{pop()}--\code{push()}: a min-heap holding the $k$ heads.
Each output step pops one head (the winner) and pushes that list's successor; the heap has $\bigO{\log k}$
not $\bigO{k}$ per emit, so $\bigO{N\log k}$ total (\figref{c07-multimerge}). This is not a special trick --
it is external sorting, $k$-way database merge, and the priority-queue schedules of the graph chapters, all
the same scheduler.

The same ``scheduler'' lens gives best-first search: to list the $k$ smallest pair sums of two sorted arrays
without materialising $n^2$ of them, treat the grid $(i,j)$ as a landscape whose value $a[i]+b[j]$ grows
rightward and downward; the smallest unexplored cell always sits on a \term{frontier}, so seed the heap at
$(0,0)$, pop the best cell, and push its two neighbours -- $k$ pops give the $k$ answers. And the median
needs the ``middle'' not an extreme, so keep \emph{two} heaps, a max-heap of the low half and a min-heap of
the high half, their sizes within one: the two roots \emph{are} the middle.

\begin{figure}[H]
\centering
\begin{tikzpicture}[diagram, x=13mm, y=8mm,
   cell/.style={draw, minimum size=6mm, inner sep=1pt, font=\small},
   flow2/.style={->, >=stealth, thick, draw=darkgray}]
  % three input lists (head highlighted)
  \node[cell, fill=primaryblue!15, draw=primaryblue] (a0) at (0,0) {$1$};
  \node[cell] (a1) at (0.95,0) {$4$}; \node[cell] (a2) at (1.9,0) {$5$};
  \node[cell, fill=primaryblue!15, draw=primaryblue] (b0) at (0,-1.1) {$1$};
  \node[cell] (b1) at (0.95,-1.1) {$3$}; \node[cell] (b2) at (1.9,-1.1) {$4$};
  \node[cell, fill=primaryblue!15, draw=primaryblue] (c0) at (0,-2.2) {$2$};
  \node[cell] (c1) at (0.95,-2.2) {$6$};
  \draw[flow2] (a0)--(a1); \draw[flow2] (a1)--(a2);
  \draw[flow2] (b0)--(b1); \draw[flow2] (b1)--(b2);
  \draw[flow2] (c0)--(c1);
  \node[annot, font=\scriptsize\bfseries, text=primaryblue] at (2.9,-1.1) {heads $\to$ heap};
  % heap box
  \node[boxnode, fill=difforange!12, draw=difforange, minimum width=17mm, minimum height=9mm]
       (heap) at (5.6,-1.1) {min-heap\\ $\{1,1,2\}$};
  \draw[flow2] (a0) .. controls (3.2,-0.2) .. (heap);
  \draw[flow2] (b0) .. controls (3.2,-1.1) .. (heap);
  \draw[flow2] (c0) .. controls (3.2,-2.0) .. (heap);
  % output
  \node[cell, fill=successgreen!18, draw=successgreen] (o) at (5.6,-2.9) {$1$};
  \draw[flow2] (heap) -- (o) node[midway, right, font=\scriptsize] {pop min,};
  \node[annot, align=left, font=\scriptsize] at (7.9,-2.9) {push its\\ successor $4$};
\end{tikzpicture}
\caption{Multiway merge of the three lists $(1,4,5),(1,3,4),(2,6)$ -- LeetCode 23's first sample. Only the
\emph{current head} of each list enters the min-heap (highlighted heads). The heap \code{pop}s the global
smallest head (here a $1$), it is appended to the output, and \emph{that list's} successor is pushed back, so
the heap always holds exactly one candidate per non-exhausted list. $N$ outputs cost $\bigO{N\log k}$.
}\label{fig:c07-multimerge}
\end{figure}

\subsection{The Multiway-Merge Skeleton}

\begin{lstlisting}[caption={k-way merge with a min-heap of list heads: pop the winner, push its successor}, label={lst:c07-multimerge}]
// k sorted vectors; emit one globally-sorted stream in O(N log k)
struct Cmp {                                // min-heap on current value
    bool operator()(const pair<int,int>& a, const pair<int,int>& b) const {
        return a.first > b.first;           // ">" => smaller value on top
    }
};
priority_queue<pair<int,int>, vector<pair<int,int>>, Cmp> pq;  // {value, listIndex}
vector<int> pos(k, 0);                      // read cursor per list
for (int i = 0; i < k; ++i)
    if (!lists[i].empty()) pq.push({lists[i][0], i});   // seed one head per list
vector<int> out;
while (!pq.empty()) {
    auto [val, i] = pq.top(); pq.pop();     // global smallest head
    out.push_back(val);
    if (++pos[i] < (int)lists[i].size())    // advance that list, refill heap
        pq.push({lists[i][pos[i]], i});
}
\end{lstlisting}

\begin{complexitybox}
The heap never holds more than $k$ items, so each of the $N$ emitted elements costs one pop and (usually)
one push at $\bigO{\log k}$: $\bigO{N\log k}$ total versus the $\bigO{Nk}$ head-scan -- the win when $k$ is
large and $N$ is bigger (external merge with $k$ sorted runs on disk). Memory $\bigO{k}$. Pairwise merging
($\bigO{N\log k}$ by merging in a tournament order) matches the bound but this head-heap version is the one
that streams, uses $\bigO{k}$ memory, and generalises unchanged to linked lists (\probref{c07-mergek}) and to
$k$ iterators.
\end{complexitybox}


\subsection{The Two-Heap Running Median}

The median is the one statistic a plain heap cannot give -- it wants the \emph{middle}, not an extreme. The
repair is to run \emph{two} heaps that meet in the middle: a \textbf{max}-heap for the low half and a
\textbf{min}-heap for the high half, kept within one element of each other. The top of each is a neighbour of
the median, so the median is read off the two roots in $\bigO{1}$ and every insert is $\bigO{\log n}$.

\begin{lstlisting}[caption={Two-heap running median: a max-heap (low) and min-heap (high) balanced to within one}, label={lst:c07-median}]
priority_queue<int> low;                                        // max-heap: lower half
priority_queue<int, vector<int>, greater<int>> high;            // min-heap: upper half

void add(int x) {
    if (low.empty() || x <= low.top()) low.push(x); else high.push(x);
    // rebalance so |low| >= |high| and differ by at most 1
    if ((int)low.size() > (int)high.size() + 1) { high.push(low.top()); low.pop(); }
    if ((int)high.size() > (int)low.size())     { low.push(high.top());  high.pop(); }
}
double median() {
    if (low.size() == high.size()) return (low.top() + high.top()) / 2.0;
    return low.top();                    // low holds the extra element when odd
}
\end{lstlisting}

\begin{notebox}[Where this lives vs.\ Chapter~8 and Chapter~9]
This is the machinery behind \emph{Find Median from Data Stream} (LeetCode 295); that problem is \emph{placed
in Chapter~8's mixed set} (\chref{c06}) so we cite the pattern here without re-listing it. Contrast two close
cousins you already own: the \emph{sliding-window maximum} (Chapter~9, \chref{c26}) uses a monotonic
\emph{deque}, not a heap, because elements expire by \emph{position} (a known index leaves the window) --
somewhere a heap would have to delete an arbitrary interior node. A \emph{sliding-window median} needs
\emph{both} tricks at once (two heaps plus lazy deletion of outgoing values): that is the harder hybrid, and
knowing which single tool fits which expiry rule is the point.
\end{notebox}

\subsection{Variations and Twists}

\begin{itemize}[nosep]
  \item \textbf{Smallest range covering $k$ lists} (LC 632): run multiway merge but keep the last element
        from \emph{each} list in the heap; the window $[\min,\max]$ of the $k$ current heads is a candidate
        range -- advance the list owning the current minimum. \hfill\lv{5}
  \item \textbf{Best-first with visited set} (LC 373, 378): the grid frontier can be reached by two parents,
        so mark cells visited (or push each cell once) or you double-count; the classic off-by-structure bug. \hfill\lv{4}--\lv{5}
  \item \textbf{Greedy with repair} (IPO, LC 502): sort projects by capital, and as your capital $w$ grows,
        move every newly-affordable project into a \emph{max}-heap of profits and take the top $k$ times --
        the heap is a ``best affordable now'' that only ever gains options. \hfill\lv{5}
  \item \textbf{$k$-way merge of iterators / files}: the \coderef{c07-multimerge} skeleton is unchanged when
        each source is a stream with a \code{next()}; this is external sort's merge phase verbatim. \hfill\lv{4}
\end{itemize}

\subsection{Common Pitfalls}

\begin{warningbox}[Trap 1 -- best-first without a visited set]
In ``$k$ smallest pairs / matrix cells'' a state can be pushed by two different parents (e.g.\ $(i{+}1,j)$ and
$(i,j{+}1)$ both reach $(i{+}1,j{+}1)$ in some schemes). Without a \code{visited} guard you re-expand and the
heap fills with duplicates -- wrong answers and blow-up. Seed only $(0,0)$, expand right/down, and mark cells
as pushed, not as popped.
\end{warningbox}
\begin{warningbox}[Trap 2 -- letting the two median heaps drift out of balance]
The two-heap median is only correct while $|\mathit{low}| - |\mathit{high}| \in \{0,1\}$ \emph{and} every
element of \code{low} $\le$ every element of \code{high}. Inserting without the rebalance steps (or forgetting
to re-check the max-heap's root against the min-heap's) silently breaks the ordering across the split, and
the ``median'' is neither half's true neighbour.
\end{warningbox}
\begin{warningbox}[Trap 3 -- heap where a deque is right (or vice-versa)]
If items leave by \emph{value} (you always remove the current extreme), a heap is correct. If they leave by
\emph{position} (a fixed-size window slides), the monotonic deque of Chapter~9 (\chref{c26}) wins because a
heap cannot cheaply remove the specific element that just aged out. Reaching for \code{priority\_queue} to
``get the max of a window'' costs an $\bigO{\log n}$ arbitrary-delete workaround that the deque gives in
$\bigO{1}$ amortised.
\end{warningbox}

\subsection{Problem Lab}

The anchor is LeetCode 23 -- merge $k$ sorted linked lists -- the purest head-heap multiway merge, only with
list nodes instead of array cursors. Once \coderef{c07-multimerge} is understood, the linked-list version is a
three-line change of ``successor'' from ``next index'' to ``\code{node->next}''.

% ---------------------------------------------------------------------------
\solvedproblem{c07-mergek}{Merge k Sorted Lists}{LeetCode 23}
  {https://leetcode.com/problems/merge-k-sorted-lists/}{LC Hard · L4}
  {min-heap of list heads, multiway merge, custom comparator on nodes, dummy-tail splicing}

\begin{problemstatement}[Merge $k$ Sorted Lists]
You are given an array of $k$ linked-list \texttt{lists}, each sorted in ascending order. Merge all the
linked lists into \textbf{one} sorted linked list and return it. Constraints: $0 \le k \le 10^{4}$, the total
number of nodes across all lists is at most $10^{4}$, and node values are in $[-10^{4}, 10^{4}]$; each list
is sorted in non-decreasing order.
\end{problemstatement}

\begin{attemptgate}[Attempt it yourself first -- \lv{4}, give yourself 25 minutes]
Do not merge lists two at a time into a heap of \emph{values}; keep a heap of \emph{list heads} (pointers), so
you never detach a node you have not yet emitted.
\peek{At any output step, the next node of the merged list must be the smallest among what set of nodes?}
\nudge{Among the \emph{current head} of each non-empty list. Put all $k$ heads in a min-heap keyed by
\code{val}; pop the smallest, append it, and push \emph{that node's} \code{next}.}
\fullhint{Use a dummy head and a running \code{tail} to splice nodes. The comparator orders \code{ListNode*}
by \code{->val} (remember \code{priority\_queue} wants ``is $a$ worse than $b$'' so use \code{a->val > b->val}
for a min-heap). Terminate when the heap is empty. For the sample $[[1,4,5],[1,3,4],[2,6]]$ the emission
order is $1,1,2,3,4,4,5,6$.}
\end{attemptgate}

\solstep{Understanding the Problem}
The merged list's next node is always the smallest of the $k$ current heads -- exactly the multiway-merge
invariant. Because the lists are linked, we can \emph{move} nodes (relink \code{next} pointers) rather than
copy values, and we never need to know a list's length. The subtlety is only that the heap stores
\emph{pointers} whose key is \code{->val}, so we supply a comparator; and we must push a popped node's
\code{next} (its successor becomes the new candidate from that list).

\solstep{Key Observations}
\begin{observation}
The heap holds at most $k$ live heads, so each of the $N$ nodes costs one pop plus one push: $\bigO{N\log k}$,
matching \coderef{c07-multimerge} and beating the $\bigO{Nk}$ head-scan.
\end{observation}
\begin{observation}
Empty lists must not enter the heap (guard \code{if (h)}) or a null deref in the comparator is \re.
\end{observation}
\begin{observation}
A \code{ListNode*} has no natural \code{operator<} by value, so the custom comparator is mandatory; the
\code{priority\_queue} convention (\code{Cmp(a,b)} true means $a$ sits below $b$) means we write
\code{a->val > b->val} to get a min-heap on values.
\end{observation}

\solstep{Think It Through With Me}
\begin{thinkbox}
\deadend{Collect all values, sort, rebuild a list.}
{$\bigO{N\log N}$ and $\bigO{N}$ extra nodes; ignores that the $k$ sortedness already gives
$\bigO{N\log k}$, and fails if lists are huge streamed inputs. Misses the whole multiway point.}
\deadend{Merge lists pairwise, repeatedly (divide and conquer).}
{Correct and $\bigO{N\log k}$ if done in a tournament order, but $\bigO{\log k}$ passes over the data and
awkward with $k$ lists of very different length; the single head-heap pass is simpler and streams.}
\aha{A min-heap of the $k$ current \emph{heads}: pop smallest, splice it out, push its \code{next}. One pass
over every node, $\bigO{\log k}$ heap ops each -- \coderef{c07-multimerge} with pointers.}
\end{thinkbox}

\solstep{Solution}
Build a min-heap of \code{ListNode*} ordered by \code{->val}. Seed it with each non-empty list's head.
Maintain a dummy node and a \code{tail} pointer. While the heap is non-empty: pop the smallest head, attach it
after \code{tail}, advance \code{tail}, and if that node has a \code{next}, push it. Return
\code{dummy.next}. This is \coderef{c07-multimerge} with the array cursor \code{pos[i]} replaced by the node's
own \code{next} pointer.

\solstep{Proof of Correctness}
\begin{proofbox}[Every node emitted in global sorted order]
Invariant: the heap contains exactly the current head of each list that still has unconsumed nodes, so its
minimum is the smallest unconsumed node over all lists. Initially each head is in the heap (seeded), true. On
a step we pop the minimum $m$ -- by the invariant no other unconsumed node is smaller, so $m$ is the correct
next element of the merged order (sortedness within its list means all its list's remaining nodes are
$\ge m$, and other lists' heads are $\ge$ the heap min). We then push $m\to$next, restoring the invariant for
that list (its new head is the next unconsumed there). Inductively every emission is the smallest remaining
node, so the output is globally non-decreasing, and every node is pushed exactly once and popped exactly once,
so all $N$ appear. \end{proofbox}

\solstep{Implementation}
\begin{lstlisting}[caption={Merge k Sorted Lists -- min-heap of live list heads (LeetCode 23)}, label={lst:c07-mergek-soln}]
class Solution {
public:
    struct Cmp {                       // min-heap on node value
        bool operator()(ListNode* a, ListNode* b) const { return a->val > b->val; }
    };
    ListNode* mergeKLists(vector<ListNode*>& lists) {
        priority_queue<ListNode*, vector<ListNode*>, Cmp> pq;
        for (ListNode* h : lists) if (h) pq.push(h);   // seed non-empty heads only
        ListNode dummy(0), *tail = &dummy;
        while (!pq.empty()) {
            ListNode* node = pq.top(); pq.pop();       // globally smallest head
            tail->next = node; tail = node;            // splice it into the output
            if (node->next) pq.push(node->next);       // that list's successor rejoins
        }
        return dummy.next;
    }
};
\end{lstlisting}

\solstep{Trace}
\begin{tracebox}[Official sample: \texttt{lists = [[1,4,5],[1,3,4],[2,6]]}]
Seed heads $\{1_A,1_B,2_C\}$. Pop $1_A$, push $4_A$ $\to\{1_B,2_C,4_A\}$. Pop $1_B$, push $3_B$
$\to\{2_C,3_B,4_A\}$. Pop $2_C$, push $6_C$ $\to\{3_B,4_A,6_C\}$. Pop $3_B$, push $4_B$
$\to\{4_A,4_B,6_C\}$. Pop $4_A$, push $5_A$ $\to\{4_B,5_A,6_C\}$. Pop $4_B$ (no next) $\to\{5_A,6_C\}$.
Pop $5_A$ $\to\{6_C\}$. Pop $6_C$ $\to\varnothing$. Output $1,1,2,3,4,4,5,6$. \checkmark
\end{tracebox}

\solstep{Complexity}
\begin{complexitybox}
Time $\bigO{N\log k}$ ($N$ total nodes, heap size $\le k$); memory $\bigO{k}$ for the heap plus the $\bigO{1}$
relinking (we move nodes, not copy). Against the naive $\bigO{Nk}$ head-scan this is the multiway-merge win,
and against full sort it is $\bigO{N\log k}$ vs.\ $\bigO{N\log N}$ -- better when $k \ll N$. Verdict: \ac.
\end{complexitybox}

\solstep{Pitfalls and Variants}
\begin{warningbox}
Ways to fall: (i) forgetting \code{if (h)} when seeding, pushing a null and then dereferencing it in the
comparator -- \re; (ii) writing the comparator as \code{a->val < b->val}, turning the min-heap into a
max-heap and emitting each list \emph{backwards}; (iii) forgetting to \code{tail->next = nullptr} at the very
end if your splicing left a dangling link (here the last node's own \code{next} is already null, so the loop
is safe -- but a manual splice in other formulations is not). \emph{Variant:} the divide-and-conquer pairwise
merge is also $\bigO{N\log k}$ with no comparator -- worth knowing because interviewers ask for it after the
heap.
\end{warningbox}

\begin{insightbox}[Key takeaways]
\begin{enumerate}[nosep, leftmargin=1.6em]
  \item \textbf{One idea, four outfits}: ``heap of the best-each-source-can-still-offer'' is multiway merge
        (\probref{c07-mergek}); ``heap of the unexplored frontier'' is best-first (LC 373/378); ``two heaps
        meeting in the middle'' is the running median (\coderef{c07-median}).
  \item \textbf{$\bigO{N\log k}$, not $\bigO{Nk}$, not $\bigO{N\log N}$}: the heap of \emph{heads} scales the
        merge to $k$ streams at logarithmic cost -- external sort's engine.
  \item \textbf{Heap vs.\ deque}: remove-by-\emph{value} $\Rightarrow$ heap; remove-by-\emph{position}
        (sliding window) $\Rightarrow$ the monotonic deque of Chapter~9 (\chref{c26}).
\end{enumerate}
\end{insightbox}

\begin{hardeningbox}
\begin{itemize}[nosep, leftmargin=1.4em]
  \item \textbf{Port \coderef{c07-multimerge}} to merge $k$ sorted arrays into a stream and verify the sample
        of \figref{c07-multimerge} yields $1,1,2,3,4,4,5,6$.
  \item \textbf{Median check}: feed $5,1,3$ into \coderef{c07-median}'s \code{add} and read \code{median()}
        after each ($5$, then $3$, then $3$) -- confirm the low/high balance holds.
  \item \textbf{Best-first sketch}: seed $(0,0)$ for LC 373 and argue why the frontier never needs the full
        $n^2$ grid -- connect it to the visited-set trap above.
\end{itemize}
\end{hardeningbox}

\subsection{Practice Problems}

\begin{stuckbox}
\begin{itemize}[nosep, leftmargin=1.4em]
  \item Did your heap hold \emph{heads / frontier cells}, not all data, so memory stayed $\bigO{k}$?
  \item For best-first: did you mark cells visited \emph{when pushed} so no cell is expanded twice?
  \item For the median: did you keep the halves ordered (\code{low} $\le$ \code{high}) and sized within $1$?
\end{itemize}
\end{stuckbox}

\begin{practiceset}
  \practice{Find K Pairs with Smallest Sums}{LeetCode 373}
    {https://leetcode.com/problems/find-k-pairs-with-smallest-sums/}{LC Medium · L4}
    {Two sorted arrays \texttt{nums1}, \texttt{nums2}; return the $k$ pairs with the smallest sums. Best-first
    on the implicit grid: a min-heap seeded with rows $(i,0)$ for $i<\min(k, n_1)$ (or just $(0,0)$), pop the
    smallest-sum cell and push its right neighbour $(i,j{+}1)$; stop after $k$ pops. Cap the heap at $k$ so you
    never store the $n_1 n_2$ products -- this is \coderef{c07-multimerge} with the ``lists'' being, for each
    $i$, the stream $\texttt{nums1[i]}+\texttt{nums2[j]}$ over $j$.}
  \practice{Kth Smallest Element in a Sorted Matrix}{LeetCode 378}
    {https://leetcode.com/problems/kth-smallest-element-in-a-sorted-matrix/}{LC Medium · L4}
    {Rows and columns sorted; find the $k$-th smallest. Same best-first frontier from $(0,0)$ expanding
    right and down with a \emph{visited} set (Trap~1), popping $k$ times; or value-binary-search counting cells
    $\le$ mid in $\bigO{n\log(\max-\min)}$. The heap version generalises the moment the ``sorted matrix'' is
    really $n$ sorted lists -- recognising that is the exercise.}
  \practice{IPO}{LeetCode 502}
    {https://leetcode.com/problems/ipo/}{LC Hard · L4}
    {Given $k$ allowed projects, each with a capital \texttt{cost[i]} and \texttt{profit[i]}, start capital $w$;
    pick up to $k$ projects (each adds its profit to $w$) to maximise final capital. \textbf{Greedy with
    repair}: sort projects by cost; as $w$ grows, move every now-affordable project into a \emph{max}-heap of
    profits and take \code{top()} for the next project. The heap is ``best affordable right now,'' refilled --
    not a merge, but the same extract-extreme-then-reinsert rhythm as \probref{c07-stone} and the Summit
    (\probref{c07-summit}).}
\end{practiceset}

\begin{recallbox}[10-second refresher -- 7.4 Classic Heap Problems]
The heap is an \emph{output scheduler}: a min-heap of \emph{list heads} merges $k$ sorted streams in
$\bigO{N\log k}$ (\coderef{c07-multimerge}, \probref{c07-mergek}); a min-heap of the \emph{frontier} answers
$k$-best pair/matrix questions without enumerating; a \emph{max}-heap of affordable options does greedy-with-
repair (IPO); and \emph{two} balanced heaps give a running median in $\bigO{\log n}$ (\coderef{c07-median}) --
LeetCode 295 itself lives in Chapter~8 (\chref{c06}). Remove-by-value wants a heap; remove-by-position (a
sliding window) wants Chapter~9's deque (\chref{c26}).
\selftest{
  \item Why does a head-heap merge $k$ lists in $\bigO{N\log k}$ while a head-scan is $\bigO{Nk}$?
  \item State the two-heap median invariant (sizes and cross-ordering) and why the median is the roots.
  \item In best-first $k$-smallest pairs, why mark cells visited when \emph{pushed} rather than popped?
  \item Give one problem each where a deque beats a heap and where a heap beats a deque.
  \item What makes IPO ``greedy with repair,'' and how does its loop echo \probref{c07-stone}?
}
\end{recallbox}
